
\chapter{Formulación del Proyecto}\label{formulacion}
\azul{Esta sección postula la introducción formal de la investigación. El texto debe guiar al lector desde un contexto global hacia el problema específico, justificando la intervención y delineando la solución propuesta. Asimismo, se sugiere introducir la terminología técnica desde un inicio, destacando las contribuciones fundamentales y proporcionando una visión general de la estructura del documento.}


\section{Propósito}
\azul{El análisis contextualiza el proyecto frente a trabajos previos afines. Se requiere evidenciar qué atributos hacen diferente a esta propuesta frente a la literatura existente, estimando una extensión cercana a las dos páginas.}

\begin{ejemplo}
   \moradoc{Aunque coexisten diversos sistemas de gestión académica, gran parte de ellos postula soluciones genéricas. Esta iniciativa discrepa de los enfoques tradicionales al diseñar una arquitectura adaptada a los flujos del Departamento de Ciencias de la Ingeniería, integrando módulos de seguimiento de tesis que las alternativas comerciales omiten.}
\end{ejemplo}

\section{Identificación del Problema}
\azul{El documento define detalladamente la anomalía o necesidad que el diseño busca mitigar. Resulta vital exponer las deficiencias de los métodos actuales y argumentar por qué la implementación de una nueva tecnología se hace indispensable.}
\section{Justificación y Aporte}
\azul{La argumentación sostiene la conveniencia del proyecto. Un planteamiento sólido responde interrogantes esenciales sobre su utilidad, los usuarios beneficiados y la naturaleza del problema práctico que se resuelve \cite[p. ~40]{hernandez2014metodologia}.}


\subsection{Utilidad del Sistema}
\azul{Esta subsección explica cómo la arquitectura propuesta mejorará los flujos de trabajo, optimizando procesos y centralizando información dispersa. Se sugiere organizar el contenido mediante los siguientes ejes:}

\subsubsection{Centralización e Integración de Datos}
\subsubsection{Optimización de Tiempos y Recursos}
\subsubsection{Disponibilidad de la Información}
\subsubsection{Generación Analítica y Estadísticas}



\subsection{Impacto y Aporte Institucional}

\azul{El análisis detalla los beneficios estructurales del sistema a largo plazo, abordando áreas como el apoyo a la toma de decisiones o la actualización de la plataforma tecnológica. Las siguientes directrices orientan la redacción:}

\subsubsection{Eficiencia Operativa Sistémica}
\subsubsection{Toma de Decisiones Estratégicas}
\subsubsection{Transparencia y Trazabilidad}
\subsubsection{Modernización Tecnológica Escalable}

\section{Marco Teórico}
\azul{El apartado agrupa los modelos, lenguajes y teorías que sustentan la construcción ingenieril. La revisión exhaustiva de literatura previene fallos de diseño y conduce al establecimiento de hipótesis o afirmaciones comprobables. Para ello, el estudiante recurre a diversas fuentes de información:}

\begin{itemize}
	\item \azul{Primarias (directas): Libros, artículos científicos de bases indexadas, repositorios de tesis y disertaciones.}
	\item \azul{Secundarias: Compilaciones, bases de datos bibliográficas y resúmenes de publicaciones especializadas.}
	\item \azul{Terciarias: Catálogos institucionales, revistas y otras publicaciones que agrupan fuentes de segunda mano.}
\end{itemize}





\begin{ejemplo}
\moradoc{El desarrollo del prototipo se rige bajo la metodología ágil Scrum, la cual sugiere una alta flexibilidad ante requisitos volátiles \cite[p. ~74]{pressman2010ingenieria}. La arquitectura se fundamentará en microservicios para garantizar la escalabilidad de las transacciones. Finalmente, la persistencia de datos aplicará la Tercera Forma Normal (3FN), lo que podría indicar una reducción significativa en la redundancia de registros.}
\end{ejemplo} 

\section{Estado del Arte}\label{alternativas}
\azul{La investigación explora las herramientas de software, hardware o metodologías que actualmente abordan desafíos análogos en la industria. El levantamiento de información debe culminar con una evaluación crítica.}

\subsection{Comparativa Frente a Soluciones Actuales}

\azul{Finalmente, se postula una comparativa directa —preferentemente mediante una matriz o tabla cruzada— que contraste las alternativas analizadas frente al sistema diseñado, evidenciando de manera objetiva el valor agregado y la innovación de la propuesta académica.}